\chapter{Lactate Dehydrogenase (LDH)}

\section*{What Is This Test?}
Lactate dehydrogenase (LDH) is a ubiquitous cytoplasmic enzyme that catalyzes the reversible interconversion of pyruvate and lactate, the final step of anaerobic glycolysis: Pyruvate + NADH + H\textsuperscript{+} $\leftrightarrow$ Lactate + NAD\textsuperscript{+}. LDH is present in virtually all cells in the body, with particularly high concentrations in red blood cells, cardiac muscle, skeletal muscle, liver, kidney, brain, and lungs. Because LDH is so widely distributed, it is a sensitive but non-specific marker of cell injury or turnover. Any process causing significant cell damage or increased cell turnover will elevate serum LDH \cite{huijgen1997clinical}. The total LDH measurement alone cannot identify the tissue source; for that, LDH isoenzyme fractionation is required. LDH has a serum half-life of approximately 10 hours in blood (LDH-4, 5) to 4--5 days (LDH-1, from heart). Following tissue injury, LDH rises later than AST or CK (peak at 3--5 days) and may remain elevated for 7--14 days, making it useful for late presentation of myocardial infarction (before troponin was available) and for monitoring conditions with ongoing cell turnover (malignancy, hemolysis). In modern practice, LDH is primarily used as: (1) a marker of hemolysis (with haptoglobin and reticulocyte count), (2) a tumor marker for lymphoma, germ cell tumors, and other malignancies (LDH is a component of prognostic indices for non-Hodgkin lymphoma---the International Prognostic Index, IPI), (3) a component of the criteria for diagnosing thrombotic thrombocytopenic purpura (TTP) and HELLP syndrome, (4) a marker of tissue damage in pulmonary embolism, and (5) a prognostic marker in COVID-19 and other critical illnesses \cite{tierney2023current,tietz2023textbook}.

\section*{Alternative Names}
LD; Lactic Dehydrogenase; Serum LDH; Total LDH; Lactate Dehydrogenase (LD); LD Total

\section*{Principle}
LDH is measured by a kinetic spectrophotometric method based on the lactate-to-pyruvate reaction: L-lactate + NAD\textsuperscript{+} $\rightarrow$ Pyruvate + NADH + H\textsuperscript{+} (catalyzed by LDH). The rate of increase in absorbance at 340 nm (as NADH is produced) is proportional to LDH activity. The IFCC standardized method uses L-lactate and NAD\textsuperscript{+} at pH 9.3--9.4 (alkaline conditions favor the lactate-to-pyruvate direction). Results are reported in IU/L at 37 degrees C. The reverse reaction (pyruvate-to-lactate) is approximately 2--3$\times$ faster and was used historically, causing inter-laboratory variability. The IFCC forward reaction (lactate-to-pyruvate) is now standard internationally \cite{panteghini2023serum}. Critical pre-analytical concern: Hemolysis causes massive, artifactual elevation of LDH because red blood cells contain LDH at approximately 150--200$\times$ the concentration found in normal serum. Even microscopic hemolysis (barely visible pink tint) can significantly elevate LDH. Hemolyzed samples must be rejected. This test is NOT performed by hematology analyzers.

\section*{History}
LDH was first crystallized from rat muscle by Straub in 1940. The clinical utility of LDH as a marker of myocardial infarction was established in the 1950s-1960s. The LDH isoenzymes were separated by electrophoresis in 1957 by Vesell and Bearn, allowing tissue-source identification. LDH was a cornerstone of MI diagnosis (along with CK-MB and AST) throughout the 1960s-1990s, used particularly for late-presenting MI because LDH peaks at 3--5 days and stays elevated for 10--14 days. High-sensitivity troponin has largely replaced LDH for cardiac indications. The role of LDH as a prognostic marker in non-Hodgkin lymphoma was established by the International Non-Hodgkin's Lymphoma Prognostic Factors Project in 1993 (IPI) \cite{international1993predictive}.

\section*{Why This Test Is Ordered}
\begin{itemize}
  \item Detection and monitoring of hemolytic anemia: elevated LDH (often 2--10$\times$ ULN) with elevated reticulocyte count and low haptoglobin is the classic triad. In intravascular hemolysis, LDH is disproportionately elevated relative to other markers \cite{barcellini2015clinical}
  \item Diagnosis of thrombotic thrombocytopenic purpura (TTP): markedly elevated LDH (often >1,000 IU/L) with thrombocytopenia, microangiopathic hemolytic anemia (schistocytes on smear), fever, neurological symptoms, and renal dysfunction. LDH elevation reflects both hemolysis and tissue ischemia
  \item Prognosis in non-Hodgkin lymphoma: LDH is one of five components of the IPI score; elevated LDH worsens prognosis. Also used in Hodgkin lymphoma (IPS score) and germ cell tumors (testicular cancer---LDH is a component of the IGCCCG risk classification)
  \item Monitoring response to lymphoma and leukemia therapy: falling LDH indicates treatment response; rising LDH suggests relapse or progression
  \item Evaluation of pulmonary embolism: elevated LDH with normal CK suggests pulmonary infarction; LDH is part of the Wells criteria in older PE risk stratification (not standard in modern PERC/Wells scoring)
  \item Assessment of COVID-19 severity: markedly elevated LDH is a marker of critical illness and predictor of mortality in COVID-19 patients
  \item Adjunctive marker in myocardial infarction: LDH-1 > LDH-2 (flipped LDH pattern) at 3--5 days post-infarction was historically diagnostic; now of historical interest only
  \item Diagnosis of HELLP syndrome in pregnancy: Hemolysis, Elevated Liver enzymes, Low Platelets---LDH >600 IU/L is one criterion
  \item Detection of malignant effusions: pleural or peritoneal fluid LDH helps differentiate exudative from transudative effusions
\end{itemize}

\section*{Is This a Primary or Secondary Test}
LDH is a secondary, adjunctive test used for specific indications (hemolysis, lymphoma prognosis, TTP, malignancy). It is not part of routine chemistry panels. It is highly sensitive but non-specific and must be interpreted in clinical context.

\section*{How This Test Is Performed}
A healthcare professional draws blood from a vein into a serum separator tube (gold-top) or lithium heparin tube (green-top). Hemolysis must be absolutely avoided---even the slightest visual hemolysis invalidates LDH results. The tourniquet should be applied briefly and gently. Blood should flow freely without excessive suction or probing. The sample is centrifuged promptly and analyzed on an automated chemistry analyzer. Results are available within 1--2 hours.

\section*{What Preparation Is Needed}
No fasting is required. The key pre-analytical requirement is avoidance of hemolysis: difficult venipuncture, small-gauge needles, excessive suction, shaking of the tube, or delayed processing all increase the risk of hemolysis and invalid LDH results. Strenuous exercise elevates LDH from muscle release, though less markedly than CK. Recent surgery, trauma, fractures, or intramuscular injections can elevate LDH.

\section*{Contraindications and Precautions}
No absolute contraindications. The critical pre-analytical concern is hemolysis: LDH concentration in RBCs is approximately 150--200$\times$ serum concentration. Even a hemoglobin concentration of 50 mg/dL (barely pink plasma) can double serum LDH. Therefore: (1) Hemolyzed samples must be rejected and the test redrawn with careful technique. (2) LDH is also found in high concentrations in platelets; serum LDH is slightly higher than plasma LDH because LDH is released from platelets during clotting. Lithium heparin plasma is acceptable and preferred by some. (3) LDH is unstable at room temperature; samples should be centrifuged and separated from cells within 2 hours. LDH-4 and LDH-5 (liver, muscle) are more labile than LDH-1 (heart) and degrade at room temperature. (4) Drugs that can elevate LDH: many chemotherapeutic agents (tumor lysis), anesthetics, aspirin, clofibrate, fluorides, mithramycin, and narcotics. (5) Macro-LDH: rare complex of LDH with immunoglobulins causing persistently elevated LDH; differentiated by PEG precipitation or electrophoresis.

\section*{Risks}
Minimal risk from venipuncture. Mild pain or bruising at site resolves in 1--2 days.

\section*{What Happens After the Test}
Results available within 1--2 hours. LDH interpretation depends heavily on clinical context: (1) LDH + low haptoglobin + elevated reticulocyte count: hemolytic anemia. (2) LDH + schistocytes + thrombocytopenia: TTP/HUS $\rightarrow$ hematology emergency requiring plasma exchange. (3) LDH + lymphadenopathy + B symptoms: lymphoma staging (IPI score). (4) LDH elevated >2$\times$ ULN in Covid-19: indicator of severe disease and predictor of mortality. (5) LDH + right upper quadrant pain + thrombocytopenia + liver enzyme elevation in third trimester: HELLP syndrome. (6) Isolated LDH elevation without obvious cause: consider macro-LDH, occult malignancy, or lab artifact from hemolysis.

\section*{Understanding Results}

\subsection*{Below Normal}
\begin{itemize}
  \item Low LDH is not clinically significant. LDH-lowering drugs (clofibrate) are not used therapeutically for this purpose.
  \item Uremia may cause falsely low LDH in vitro (urea inhibits LDH activity).
  \item High-dose ascorbic acid (vitamin C) can cause falsely low LDH by interfering with the NADH detection method.
\end{itemize}

\subsection*{Above Normal}
\begin{itemize}
  \item \textbf{Hemolytic anemia (most specific use):} Intravascular hemolysis (autoimmune hemolytic anemia, transfusion reaction, PNH, mechanical heart valve, microangiopathic hemolytic anemia---TTP, HUS, DIC, HELLP). Extravascular hemolysis (hereditary spherocytosis, G6PD deficiency, sickle cell disease, thalassemia---LDH elevated but typically less so than intravascular). Hemolysis labs: elevated LDH, elevated indirect bilirubin, low haptoglobin, elevated reticulocyte count. Coombs test positive in autoimmune hemolytic anemia.
  \item \textbf{Thrombotic thrombocytopenic purpura (TTP):} Markedly elevated LDH (often >1,000--2,000 IU/L) reflecting both hemolysis and tissue ischemia from microvascular thrombosis. ADAMTS13 activity <10\% is diagnostic. Pentad: thrombocytopenia, microangiopathic hemolytic anemia, neurological symptoms, renal dysfunction, fever (complete pentad present in only 5--40\% of cases). Plasma exchange is emergent treatment---TTP is fatal without intervention \cite{george2014syndromes}.
  \item \textbf{Malignancy (broadest association):} (a) Lymphomas: Hodgkin and non-Hodgkin lymphoma. LDH is one of five IPI factors (age >60, performance status >1, advanced stage, extranodal sites >1, elevated LDH). Elevated LDH independently predicts poor prognosis and is used for risk stratification. (b) Leukemia: ALL, AML, CML blast crisis. (c) Germ cell tumors: testicular cancer (LDH is a tumor marker and one of three IGCCCG risk factors along with AFP and beta-hCG). (d) Metastatic solid tumors: liver, breast, lung, melanoma. (e) Tumor lysis syndrome: LDH rises with massive cell lysis after chemotherapy.
  \item \textbf{Myocardial infarction (historical):} LDH rises 24--72 hours after MI, peaks at 3--5 days, normalizes at 10--14 days. LDH-1 > LDH-2 (flipped LDH pattern) was diagnostic of MI. Now of historical interest only; troponin is the standard biomarker.
  \item \textbf{Pulmonary embolism/infarction:} Elevated LDH with normal CK and elevated bilirubin suggests pulmonary infarction (tissue necrosis). LDH >600 IU/L is a minor criterion in older PE diagnostic scoring systems.
  \item \textbf{Liver disease:} Hepatitis, cirrhosis, hepatic congestion (right heart failure, Budd-Chiari). LDH is elevated but ALT is usually proportionally higher. In ischemic hepatitis (shock liver), ALT and LDH are both markedly elevated but ALT rise is more dramatic.
  \item \textbf{Other:} Skeletal muscle injury (rhabdomyolysis, myositis), acute pancreatitis, renal infarction, seizures, prolonged hypotension/shock, COVID-19, infectious mononucleosis (EBV), megaloblastic anemia (ineffective erythropoiesis).
\end{itemize}

\section*{Normal Results}
\begin{itemize}
  \item \textbf{Adults:} 140--280 IU/L (varies by method and laboratory; typical 100--200 IU/L) \cite{wallach2022interpretation}
  \item \textbf{Children:} 150--400 IU/L (higher in children due to growth and higher cell turnover)
  \item \textbf{Newborns:} 160--450 IU/L (highest in first week of life)
  \item \textbf{LDH half-lives:} LDH-1 (heart): ~4--5 days; LDH-5 (liver): ~10 hours
  \item \textbf{LDH peak post-MI:} 3--5 days; normalization 10--14 days
  \item \textbf{LDH in hemolysis:} Often 2--10$\times$ ULN
\end{itemize}

\section*{Follow-up Tests}
\begin{itemize}
  \item \textbf{LDH isoenzyme fractionation:} Identifies tissue source when origin is unclear.
  \item \textbf{Haptoglobin, reticulocyte count, CBC, peripheral smear:} For hemolytic anemia workup.
  \item \textbf{Direct Coombs test:} For autoimmune hemolytic anemia.
  \item \textbf{ADAMTS13 activity:} For TTP diagnosis.
  \item \textbf{Imaging (CT, PET-CT, ultrasound):} For lymphoma staging or suspected malignancy.
  \item \textbf{Liver function tests (CMP):} AST/ALT/ALP/bilirubin to evaluate liver involvement.
  \item \textbf{Tumor markers:} AFP, beta-hCG (germ cell tumors). LDH is itself a tumor marker.
  \item \textbf{Cardiac troponin:} For MI diagnosis (replaces LDH).
  \item \textbf{CT pulmonary angiogram:} For suspected PE.
\end{itemize}

\section*{References}
\bibliographystyle{unsrtnat}
\bibliography{references}

\clearpage
